\subsection{辫群与映射类群}\label{sec:20f}
现在专门考察最简单的非形式情形 $n=1$、$m=5$，并换一个视角来理解其中的情形，从而说明高阶复合映射的几何意义。将等价\eref{20.3}与 Yoneda 嵌入结合，得到全忠实 $\Ai$ 函子
\begin{equation}\label{eq:20.18}
\F(X)\longrightarrow\Q=\modcat(\F).
\end{equation}
它把正合 Lagrangian 膜 $L^\#$ 映到 $\Ai$ 模 $\M$，其上同调为
\[
M=H(\M)=\bigoplus_i HF^*(V_i^\#,L^\#),
\]
将其视为 $A=\bigoplus_{i,j}HF^*(V_i^\#,V_j^\#)$ 上的模。我们特别关心\sr{fig:20.2}{图~20.2}所示的曲线 $L$。它的 Maslov 数为零，并且可使它正合，因此它成为一个正合 Lagrangian 膜。适当选择分次后，对应的模为
\[
M=U_3\oplus SU_3,
\]
其中 $U_i=(i)R\cong\K$ 表示初等单模，即将 $R$ 视为右 $A$ 模时的各直和项；$S$ 像往常一样表示平移。
\SourceFigure{figure-20-2}{0.458999506}{20.2}
接下来更仔细地考察单模 $U_3$ 的同调代数。以 $P_i=(i)A$ 表示初等投射模。存在无上同调复形
\[
\begin{aligned}
0\longrightarrow S^{-3}U_3\longrightarrow S^{-2}P_3
&\longrightarrow S^{-1}P_2\oplus S^{-2}P_4\\
&\longrightarrow P_1\oplus S^{-1}P_3\oplus S^{-2}P_5\\
&\longrightarrow P_2\oplus S^{-1}P_4\longrightarrow P_3\longrightarrow U_3\longrightarrow0,
\end{aligned}
\]
\SourcePage{321}
可将它首尾接续，得到 $U_3$ 的无限投射分解。由此计算得到下列结果；其中的分次记号与 Hochschild 上同调群的记号类似：
\begin{equation}\label{eq:20.19}
\begin{aligned}
&\operatorname{Ext}_{A}^{r+s}(U_3,U_3)^s=0,
&&r+s<0;\\
&\operatorname{Ext}_{A}^{0}(U_3,U_3)^0=\K,\qquad
\operatorname{Ext}_{A}^{0}(U_3,U_3)^s=0,
&&s\ne0;\\
&\operatorname{Ext}_{A}^{1}(U_3,U_3)^{-1}=\K,\qquad
\operatorname{Ext}_{A}^{1}(U_3,U_3)^s=0,
&&s\ne-1.
\end{aligned}
\end{equation}
最后，在已知这些结果的基础上，使用条形分解，或者与扩充代数 $\overline A$ 的对应结果比较，还可看出乘积
\begin{equation}\label{eq:20.20}
\operatorname{Ext}_{A}^{1}(U_3,U_3)^{-1}\otimes
\operatorname{Ext}_{A}^{1}(U_3,U_3)^{-1}
\longrightarrow\operatorname{Ext}_{A}^{2}(U_3,U_3)^{-2}
\end{equation}
非零。这些计算的意义如下。令 $\widetilde{\A}$ 为 $A$ 上的平凡 $\Ai$ 结构，并令 $\widetilde{\Q}=\modcat(\widetilde{\A})$ 为其 $\Ai$ 模范畴。对于上同调模 $M$ 固定的 $\widetilde{\A}$ 模 $\widetilde{\M}$，有一种一般的分类方法，在 Kodaira--Spencer 形变理论的意义下由 $\operatorname{Ext}_{A}(M,M)$ 控制。在此背景下，\eref{20.19}和\eref{20.20}描述了一种一阶形变空间为一维、且所有这些一阶形变在二阶均受阻的理论。\znote{313}{把 (20.19)(20.20) 读成形变论：$\operatorname{Ext}^1_A(U_3,U_3)^{-1}=\K$ 是一维的一阶形变空间，(20.20) 非零则说明两个一阶形变之积给出非零的二阶障碍，即每个一阶形变都不可积。于是 $U_3$ 上只存在平凡的 $\widetilde{\A}$ 模结构，这正是\stref{引理}{20.6}的出发点。符号的配位也相容：形变理论中一阶形变住在 $\operatorname{Ext}^1$，障碍住在 $\operatorname{Ext}^2$，而此处内部分次分别为 $-1$ 与 $-2$。[来源：A，原书本节。]}因此，在 $H^0(\widetilde{\Q})$ 的同构意义下，将 $U_3$ 赋予 $\widetilde{\A}$ 模结构 $\widetilde U_3$ 的方式只有一种，即平凡的方式。使用同一形式体系，通过略为复杂的论证，还可得到下列结论；这里略去其细节。
\StatementNumber{6}\begin{lemma}\label{st:20.6}
设 $\widetilde{\M}$ 为 $\widetilde{\A}$ 上的 $\Ai$ 模，其底层上同调模为 $H(\widetilde{\M})\cong U_3\oplus SU_3$。那么，在 $H^0(\widetilde{\Q})$ 中有
\[
\widetilde{\M}\cong\ACone(t),
\]
其中 $t\in\hom_{\widetilde{\Q}}^0(\ModShift^{-1}\widetilde U_3,\ModShift\widetilde U_3)$ 表示 $\operatorname{Ext}_{A}^{2}(U_3,U_3)$ 中的一个元素。因此，
\[
H^{-1}\bigl(\hom_{\widetilde{\Q}}(\widetilde{\M},\widetilde{\M})\bigr)\ne0.
\]
\end{lemma}

回到几何，暂且假设 $\F$ 是形式的，此时 $\Q$ 与 $\widetilde{\Q}$ 拟等价。于是，\sr{fig:20.2}{图~20.2}中的曲线 $L$ 会给出一个适用\stref{引理}{20.6}的 $\Ai$ 模。另一方面，由于\eref{20.18}是忠实的，这个模的自同态应恢复 $HF^*(L^\#,L^\#)\cong H^*(L;\K)$，而它在 $-1$ 次为零，产生矛盾。这给出了\stref{推论}{20.4}的另一种证明；这里证明了 $m=5$ 的情形，一般情形可直接由此推出。

$\Q$ 和 $\widetilde{\Q}$ 都带有辫群 $\operatorname{Br}_6$ 的作用，该作用由对应于基本球面对象的扭转函子给出。不过，在几何情形 $\Q\cong\modcat(\F)$ 中，这一作用可嵌入一个更大的群的作用，即模去同痕的 $X$ 的分次
\SourcePage{322}
正合辛自同构群。沿 $L$ 的 Dehn 扭转就是这个更大的群中不属于 $\operatorname{Br}_6$ 的一个元素；这是因为 $L$ 本身的同痕类在 $\pi:X\to\CC$ 的覆盖对合下不保持不变。直观地说，可以认为 $\F$ 的非形式性与这个更大的对称群相联系。
\clearpage
\StatementNumber{7}\begin{note}\label{st:20.7}
上面的讨论中出现的辫群与映射类群之间的关系，是经典低维拓扑中已得到充分理解的一个课题\bcite{23}。较近的研究表明，广义辫群的情形更为复杂\bcite{142}。值得从 Fukaya 范畴的角度简要考察这一问题。

考虑 $(E_6)$ 奇点的 Milnor 纤维，即 $\{x_0^3+x_1^4=1\}\subset\CC^2$。在拓扑上，它是具有一个穿孔的三亏格曲面，并带有构成 $(E_6)$ 构型的一组消失圈基。将它们记为 $(V_1,\ldots,V_6)$，编号采用\bcite{142}，图~2 所示的次序，尽管这不是它们在该基中出现的自然次序。令 $\A$ 为该 Milnor 纤维的 Fukaya 范畴中以 $V_i^\#$ 为对象的全 $\Ai$ 子范畴。在上同调层次，$A=H(\A)$ 可表示为适当分次的 $(E_6)$ Dynkin 箭图的路径代数的商，其含义与之前相同。令
\[
L=\tau_{V_4}\tau_{V_3}\tau_{V_5}\tau_{V_4}\tau_{V_2}\tau_{V_6}\tau_{V_3}\tau_{V_5}\tau_{V_4}(V_1).
\]
经过适当同痕后，$L$ 与 $V_1$ 恰有一个交点，由此得到两条曲线相应的 Dehn 扭转之间的辫关系。如\bcite{142}所指出，$(E_6)$ 型广义辫群中不存在对应的关系，因此，从该辫群到 Milnor 纤维的映射类群的自然同态不是单射。\znote{314}{这里的机制是：$L$ 由这一串 Dehn 扭转作用于 $V_1$ 得到，适当同痕后与 $V_1$ 恰交一点，于是 $\tau_L\tau_{V_1}$ 与 $\tau_{V_1}\tau_L$ 之间出现了通常的辫关系；而\bcite{142}指出 $(E_6)$ 型广义辫群中没有对应关系，故该同态不单射。后半段的计算把这一失效搬到 $\Tw\widetilde{\A}$ 上：若 $\A$ 形式，则 $H(\hom(Y_1,C))$ 应与 $HF^*(V_1,L)$ 同为一维，而实算为三维（次数 $-2,-1,0$）。[来源：A，原书\stref{笔记}{20.7}与所引文献。]}

令 $\widetilde{\A}$ 为 $A$ 上的平凡 $\Ai$ 结构。在 $\Tw\widetilde{\A}$ 中考虑扭曲复形
\[
C=T_{Y_4}T_{Y_3}T_{Y_5}T_{Y_4}T_{Y_2}T_{Y_6}T_{Y_3}T_{Y_5}T_{Y_4}(Y_1),
\]
其中 $Y_i$ 为 $\widetilde{\A}$ 本身的对象。直接计算可知，$H\bigl(\hom_{\Tw\widetilde{\A}}(Y_1,C)\bigr)$ 为三维，生成元的次数分别为 $-2,-1,0$。

若 $\A$ 拟同构于 $\widetilde{\A}$，则可用\stref{推论}{17.17}将这个 $H(\hom)$ 群与 $(V_1,L)$ 的 Floer 上同调等同。但我们知道，后者是一维的，因此 $\A$ 不可能是形式的。由此可见，$\A$ 中的高阶乘积会导致扭曲复形之间的同态发生抵消，而这些抵消似乎又是 $(E_6)$ 辫群作用不忠实的原因。我们说“似乎”，是因为尚不知道该群是否忠实地作用于 $\Tw\widetilde{\A}$。
\end{note}
